\chapter{Localization of a product ring at a single-coordinate element}
\label{chap:Mathlib_RingTheory_Localization_Away_Pi}

This chapter tracks the helper
\lean{Localization.awayPiSingleEquiv}
in \texttt{Proetale/Mathlib/RingTheory/Localization/Away/Pi.lean}.

\begin{theorem}
  \label{thm:localization-away-pi-single-equiv}
  \lean{Localization.awayPiSingleEquiv}
  \leanok
  Let $S$ be a commutative ring, $I$ a type, and $\{k_i\}_{i \in I}$ a
  family of commutative $S$-algebras. Fix $i_0 \in I$ and an element
  $s \in k_{i_0}$. Then localizing the product $\prod_{i \in I} k_i$ at
  the singleton-supported element $\operatorname{single}_{i_0}(s)$
  collapses to localizing the $i_0$ factor at $s$: there is a canonical
  $S$-algebra isomorphism
  \[
    \Bigl(\prod_{i \in I} k_i\Bigr)\!\left[\tfrac{1}{\operatorname{single}_{i_0}(s)}\right]
    \;\xrightarrow{\sim}\; k_{i_0}[1/s].
  \]
\end{theorem}
\begin{proof}
  \leanok
  Write $R := \prod_{i \in I} k_i$, $r := \operatorname{single}_{i_0}(s)$,
  and let $\pi : R \to k_{i_0}$ be the projection $x \mapsto x_{i_0}$,
  viewed as an $S$-algebra homomorphism. The composite
  $f := (\text{alg. map}) \circ \pi : R \to k_{i_0}[1/s]$ sends $r$ to
  $\operatorname{algMap}(s)$, which is a unit in $k_{i_0}[1/s]$, so $f$
  sends every power of $r$ to a unit. The universal property of
  $R[1/r]$ then yields a unique $S$-algebra homomorphism
  $\widetilde{f} : R[1/r] \to k_{i_0}[1/s]$ such that
  $\widetilde{f}(x/1) = f(x)$.

  We claim $\widetilde{f}$ is bijective.

  \emph{Surjectivity.} Every element of $k_{i_0}[1/s]$ has the form
  $\operatorname{mk}'(x, s^n)$ for some $x \in k_{i_0}$ and $n \geq 0$.
  The element $\operatorname{mk}'(\operatorname{single}_{i_0}(x),\,r^n)$
  in $R[1/r]$ is a preimage, since $\pi(\operatorname{single}_{i_0}(x))
  = x$ and $\pi(r^n) = s^n$.

  \emph{Injectivity.} Suppose $\widetilde{f}(\operatorname{mk}'(x, r^n))
  = 0$. Clearing denominators gives $f(x) = \operatorname{algMap}(x_{i_0})
  = 0$ in $k_{i_0}[1/s]$. By the standard characterisation of the kernel
  of $k_{i_0} \to k_{i_0}[1/s]$ there is $m \geq 0$ with
  $s^m \cdot x_{i_0} = 0$ in $k_{i_0}$. We show
  $\operatorname{mk}'(x, r^n) = 0$ in $R[1/r]$ by exhibiting
  $r^{m+1}$ as an annihilator of $x$. Indeed, for $j \in I$:
  \begin{itemize}
    \item if $j = i_0$, $(r^{m+1} \cdot x)_j = s^{m+1} \cdot x_{i_0} =
      s \cdot (s^m \cdot x_{i_0}) = 0$;
    \item if $j \neq i_0$, $r_j = 0$, so $(r^{m+1})_j = 0^{m+1} = 0$ and
      $(r^{m+1} \cdot x)_j = 0$.
  \end{itemize}
  Hence $r^{m+1} \cdot x = 0$, so $\operatorname{mk}'(x, r^n) = 0$ in
  $R[1/r]$.

  Combining the two, $\widetilde{f}$ is an $S$-algebra isomorphism
  $R[1/r] \xrightarrow{\sim} k_{i_0}[1/s]$.
\end{proof}
